% GENERATED FILE. Edit canonical lessons, then run npm run generate:books.
% canonical-source-hash: fnv1a64:77f56e8637f96e70
% canonical-lessons: PA-C212-nahin-return, PA-C212-nahin-pinda, PA-C212-nahin-review, PA-C212-not-open, PA-C212-payoff

\chapter{I Don't Drink Tea}
\label{ch:pa-a2-negate}
\hlchaptermodality{212}

\begin{chapteropening}
\textbf{By the end of this chapter:} \emph{I can say that I do not do something, and that something is not so, with nahin before the verb.}

Say what you do not eat, drink or do, and that a place is not open, with nahin before the verb.
\end{chapteropening}

\section[\texorpdfstring{\pa{ਨਹੀਂ} · \pa{ਮੈਨੂੰ} \pa{ਸਮਝ} \pa{ਨਹੀਂ} … (nahīṁ · mainūṁ …)}{ਨਹੀਂ · ਮੈਨੂੰ ਸਮਝ ਨਹੀਂ … (nahīṁ · mainūṁ …)}]{Return to a familiar no}

\label{lesson:PA-C212-nahin-return}



\begin{quote}
\emph{Hear:} \emph{nahīṁ} Say \textbf{\pa{ਨਹੀਂ}}. Then the whole sentence you have used since the early lessons: \textbf{\pa{ਮੈਨੂੰ} \pa{ਸਮਝ} \pa{ਨਹੀਂ} \pa{ਆਈ}}, \textquotedblleft{}I did not understand\textquotedblright{}.
\end{quote}

\begin{grammarlens}[title={Grammar Lens}]
You learned that \textbf{\pa{ਨਹੀਂ}} goes directly before the verb and that nothing else in the sentence moves. Hold that idea next to a sentence you can say now: \textbf{\pa{ਮੈਂ} \pa{ਚਾਹ} \pa{ਪੀਂਦਾ} \pa{ਹਾਂ}}.
\end{grammarlens}

\subsection*{Your turn}
\begin{itemize}
\raggedright
  \item \emph{Hear:} \emph{no} Say \textbf{\pa{ਨਹੀਂ}}.
  \item \emph{Hear:} \emph{I did not understand} Say \textbf{\pa{ਮੈਨੂੰ} \pa{ਸਮਝ} \pa{ਨਹੀਂ} \pa{ਆਈ}}.
  \item \emph{Hear:} \emph{I drink tea, as a man} Say \textbf{\pa{ਮੈਂ} \pa{ਚਾਹ} \pa{ਪੀਂਦਾ} \pa{ਹਾਂ}}.
\end{itemize}

\begin{tcolorbox}[breakable,colback=teal!4,colframe=teal!35!black,title={Before you move on}]
Say \emph{no}, then \emph{I drink tea}. The next lesson joins them.
\end{tcolorbox}
\section[\texorpdfstring{\pa{ਮੈਂ} \pa{ਚਾਹ} \pa{ਨਹੀਂ} \pa{ਪੀਂਦਾ} \pa{ਹਾਂ।} (maiṁ chāh nahīṁ …)}{ਮੈਂ ਚਾਹ ਨਹੀਂ ਪੀਂਦਾ ਹਾਂ। (maiṁ chāh nahīṁ …)}]{\pa{ਮੈਂ} \pa{ਚਾਹ} \pa{ਨਹੀਂ} \pa{ਪੀਂਦਾ} \pa{ਹਾਂ} — I don't drink tea}

\label{lesson:PA-C212-nahin-pinda}



\begin{quote}
Say \emph{I drink tea every day}, with the ending that fits you.
\end{quote}

\begin{grammarlens}[title={Grammar Lens: \pa{ਨਹੀਂ} before the verb}]
Put \textbf{\pa{ਨਹੀਂ}} directly before \textbf{\pa{ਪੀਂਦਾ}}. Everything else stays where it was, the helper \textbf{\pa{ਹਾਂ}} included:

\begin{quote}
\textbf{\pa{ਮੈਂ} \pa{ਚਾਹ} \pa{ਨਹੀਂ} \pa{ਪੀਂਦਾ} \pa{ਹਾਂ।}} — \emph{maiṁ chāh nahīṁ pīndā hāṁ.} — I don't drink tea (a man speaking).
\end{quote}

\begin{quote}
\textbf{\pa{ਮੈਂ} \pa{ਚਾਹ} \pa{ਨਹੀਂ} \pa{ਪੀਂਦੀ} \pa{ਹਾਂ।}} — \emph{maiṁ chāh nahīṁ pīndī hāṁ.} — I don't drink tea (a woman speaking).
\end{quote}

\begin{quote}
\textbf{\pa{ਉਹ} \pa{ਰੋਟੀ} \pa{ਨਹੀਂ} \pa{ਖਾਂਦਾ} \pa{ਹੈ।}} — \emph{uh roṭī nahīṁ khāndā hai.} — He doesn't eat roti.
\end{quote}

This book keeps the helper after \textbf{\pa{ਨਹੀਂ}} in every sentence it prints.
\end{grammarlens}

\subsection*{Your turn}
\begin{itemize}
\raggedright
  \item \emph{Hear:} \emph{I don't drink milk, as a woman} Say \textbf{\pa{ਮੈਂ} \pa{ਦੁੱਧ} \pa{ਨਹੀਂ} \pa{ਪੀਂਦੀ} \pa{ਹਾਂ।}}
  \item \emph{Hear:} \emph{he doesn't go to school} Say \textbf{\pa{ਉਹ} \pa{ਸਕੂਲ} \pa{ਨਹੀਂ} \pa{ਜਾਂਦਾ} \pa{ਹੈ।}}
  \item \emph{Read it:} \textbf{\pa{ਮੈਂ} \pa{ਚਾਹ} \pa{ਨਹੀਂ} \pa{ਪੀਂਦਾ} \pa{ਹਾਂ।}} Point to \textbf{\pa{ਨਹੀਂ}}, then to the verb after it.
  \item \emph{Write it:} copy \textbf{\pa{ਨਹੀਂ} \pa{ਪੀਂਦਾ} \pa{ਹਾਂ}} once
\end{itemize}

\begin{tcolorbox}[breakable,colback=teal!4,colframe=teal!35!black,title={Before you move on}]
Where does \textbf{\pa{ਨਹੀਂ}} go? (Directly before the verb.) Does \textbf{\pa{ਹਾਂ}} stay? (Yes.)
\end{tcolorbox}
\section[\texorpdfstring{\pa{ਨਹੀਂ} \pa{ਪੀਂਦਾ} \pa{ਹਾਂ} (nahīṁ pīndā hāṁ)}{ਨਹੀਂ ਪੀਂਦਾ ਹਾਂ (nahīṁ pīndā hāṁ)}]{Say what you don't do}

\label{lesson:PA-C212-nahin-review}



\begin{quote}
\emph{Hear:} \emph{I don't drink tea} Say it with the ending that fits you.
\end{quote}

\subsection*{Your turn}
\begin{itemize}
\raggedright
  \item \emph{Hear:} \emph{she doesn't eat roti} Say \textbf{\pa{ਉਹ} \pa{ਰੋਟੀ} \pa{ਨਹੀਂ} \pa{ਖਾਂਦੀ} \pa{ਹੈ।}}
  \item \emph{Hear:} \emph{I don't go home, as a man} Say \textbf{\pa{ਮੈਂ} \pa{ਘਰ} \pa{ਨਹੀਂ} \pa{ਜਾਂਦਾ} \pa{ਹਾਂ।}}
  \item \emph{Read it:} \textbf{\pa{ਉਹ} \pa{ਸਕੂਲ} \pa{ਨਹੀਂ} \pa{ਜਾਂਦਾ} \pa{ਹੈ।}} Give its meaning.
  \item \emph{Write it:} cover the model and write \textbf{\pa{ਮੈਂ} \pa{ਚਾਹ} \pa{ਨਹੀਂ} \pa{ਪੀਂਦੀ} \pa{ਹਾਂ।}} once
\end{itemize}

\begin{tcolorbox}[breakable,colback=teal!4,colframe=teal!35!black,title={Before you move on}]
Say one thing you do every day, then one thing you don't.
\end{tcolorbox}
\section[\texorpdfstring{\pa{ਬਾਗ਼} \pa{ਖੁੱਲ੍ਹਾ} \pa{ਨਹੀਂ} \pa{ਹੈ।} (bāgh khullhā …)}{ਬਾਗ਼ ਖੁੱਲ੍ਹਾ ਨਹੀਂ ਹੈ। (bāgh khullhā …)}]{\pa{ਬਾਗ਼} \pa{ਖੁੱਲ੍ਹਾ} \pa{ਨਹੀਂ} \pa{ਹੈ} — the garden is not open}

\label{lesson:PA-C212-not-open}



\begin{quote}
Say \emph{the garden is open in the morning}, then \emph{the garden is closed}. Then \emph{I don't drink tea}.
\end{quote}

\begin{grammarlens}[title={Grammar Lens: when \pa{ਹੈ} is the verb}]
In \textbf{\pa{ਬਾਗ਼} \pa{ਖੁੱਲ੍ਹਾ} \pa{ਹੈ}} the verb is \textbf{\pa{ਹੈ}} itself. So \textbf{\pa{ਨਹੀਂ}} goes directly before \textbf{\pa{ਹੈ}}:

\begin{quote}
\textbf{\pa{ਬਾਗ਼} \pa{ਖੁੱਲ੍ਹਾ} \pa{ਨਹੀਂ} \pa{ਹੈ।}} — \emph{bāgh khullhā nahīṁ hai.} — The garden is not open.
\end{quote}

\begin{quote}
\textbf{\pa{ਬਾਗ਼} \pa{ਬੰਦ} \pa{ਨਹੀਂ} \pa{ਹੈ।}} — \emph{bāgh band nahīṁ hai.} — The garden is not closed.
\end{quote}

\begin{quote}
\textbf{\pa{ਉਹ} \pa{ਘਰ} \pa{ਵਿੱਚ} \pa{ਨਹੀਂ} \pa{ਹੈ।}} — \emph{uh ghar vichch nahīṁ hai.} — He is not at home.
\end{quote}
\end{grammarlens}

\subsection*{Your turn}
\begin{itemize}
\raggedright
  \item \emph{Hear:} \emph{the garden is not closed} Say \textbf{\pa{ਬਾਗ਼} \pa{ਬੰਦ} \pa{ਨਹੀਂ} \pa{ਹੈ।}}
  \item \emph{Hear:} \emph{she is not at home} Say \textbf{\pa{ਉਹ} \pa{ਘਰ} \pa{ਵਿੱਚ} \pa{ਨਹੀਂ} \pa{ਹੈ।}}
  \item \emph{Read it:} \textbf{\pa{ਬਾਗ਼} \pa{ਖੁੱਲ੍ਹਾ} \pa{ਨਹੀਂ} \pa{ਹੈ।}} Point to the word \textbf{\pa{ਨਹੀਂ}} stands before.
  \item \emph{Write it:} copy \textbf{\pa{ਨਹੀਂ} \pa{ਹੈ}} once
\end{itemize}

\begin{tcolorbox}[breakable,colback=teal!4,colframe=teal!35!black,title={Before you move on}]
\emph{Not open}? (\textbf{\pa{ਖੁੱਲ੍ਹਾ} \pa{ਨਹੀਂ} \pa{ਹੈ}}.) \emph{Doesn't drink}? (\textbf{\pa{ਨਹੀਂ} \pa{ਪੀਂਦਾ} \pa{ਹੈ}}.)
\end{tcolorbox}
\section[\texorpdfstring{\pa{ਮੈਂ} \pa{ਚਾਹ} \pa{ਨਹੀਂ} \pa{ਪੀਂਦਾ} \pa{ਹਾਂ।} (maiṁ chāh nahīṁ …)}{ਮੈਂ ਚਾਹ ਨਹੀਂ ਪੀਂਦਾ ਹਾਂ। (maiṁ chāh nahīṁ …)}]{Saying no: a four-skill payoff}

\label{lesson:PA-C212-payoff}



\begin{quote}
Say \emph{I don't drink tea}, then \emph{the garden is not open}.
\end{quote}

\subsection*{Your turn: four skills}
\begin{itemize}
\raggedright
  \item \emph{Hear:} \emph{I don't eat roti} Say it with the ending that fits you.
  \item \emph{Hear:} \emph{the garden is not closed} Say \textbf{\pa{ਬਾਗ਼} \pa{ਬੰਦ} \pa{ਨਹੀਂ} \pa{ਹੈ।}}
  \item \emph{Read it:} \textbf{\pa{ਉਹ} \pa{ਸਕੂਲ} \pa{ਨਹੀਂ} \pa{ਜਾਂਦੀ} \pa{ਹੈ।}} Give its meaning: she doesn't go to school.
  \item \emph{Write it:} cover the line and write \textbf{\pa{ਮੈਂ} \pa{ਦੁੱਧ} \pa{ਨਹੀਂ} \pa{ਪੀਂਦਾ} \pa{ਹਾਂ।}} or its woman's form
\end{itemize}

\begin{tcolorbox}[breakable,colback=teal!4,colframe=teal!35!black,title={Before you move on}]
Say two true things you don't do, and one place that is not open now.
\end{tcolorbox}
